\chapter{Verwandte Arbeiten}

Es gibt verschiedene Paper, die in Teilen Performanceunterschiede zwischen FPGAs und
Simulationen untersuchen, jedoch diese Vergleiche nicht als zentralen Inhalt betrachten.
\cite{schkufza2019just} stellen in ihrem Paper einen neuen Design-Workflow/Entwicklungsprozess
”Cascade” für FPGAs vor, in dem zunächst der Verilog-Code mittels einer Softwaresimulation ausgeführt wird,
sodass nicht auf die verschiedenen Schritte bei der Synthese eines FPGA gewartet werden muss.
Diese Schritte sowie die Übertragung des Bitstreams laufen dann im Hintergrund ab, für den*die Anwender*in ändert sich nur die
Geschwindigkeit, sobald der Synthese-Prozess abgeschlossen ist und von der Simulation auf den FPGA gewechselt wird.
In diesem Rahmen wird auch die Performance mit verschiedenen Simulatoren (auch Verilator) verglichen:
Quartus und Vivado erreichen laut Autor*innen Geschwindigkeiten von 10–100 Hz, iVerilog ca. 1 KHz und Verilator ca. 10 KHz.
Dies ist auch ein Hinweis darauf, dass Verilator tatsächlich zu den schnellsten bekannten Simulationstools gehört und Verbreitung findet.
Allerdings werden hier lediglich mögliche Taktraten verglichen und kein konkretes Benchmark-Tool genutzt.
Interessant wären auch hier Performancevergleiche mittels Benchmarks zwischen den einzelnen Simulatoren.
In Zusammenhang mit den Ergebnissen der vorliegenden Arbeit würde dies bedeuten, dass Nutzer*innen vermutlich einen Speedup um den Faktor 3.8 bemerken würden, wenn sie einen vergleichbaren Designprozess wie in Cascade beschrieben durchliefen.

Die angestellten Vergleiche zwischen Simulationssoftware und Hardware ergänzen Arbeiten, die die Geschwindigkeiten von Simulationssoftware untereinander \cite{chen2025gsim, choi2020flash, hannemann2025verilator, lopez2023fast, zhou2023khronos}
und Arbeiten, die die Geschwindigkeit von Verilator mit derer neu entwickelter Hardware vergleicht \cite{emami2024highly, qin2025feriver}.





Eine genauere Erkenntnis darüber, wie groß erwartbare Differenzen sind, kann dabei helfen, Geschwindigkeitsunterschiede zu neuer Hard- und Software besser einzuordnen.
Ein Paper, welches einen quantitativen Performancevergleich zwischen einer Verilator Simulation und einer FPGA-Synthese
für einen identischen CPU-Kern anstellt, konnte im Rahmen der Literaturrecherche nicht gefunden werden.

% https://dl.acm.org/doi/10.1145/3676641.3716010
% https://ieeexplore.ieee.org/document/4068926
% https://ieeexplore.ieee.org/document/7551388